\section{Introduction}
Low-rank training methods such as GaLore, OSD, SubTrack++, and related projection optimizers reduce optimizer state by representing gradients in a learned subspace rather than maintaining full Adam moments for every matrix entry \cite{zhao2024galore,su2025galore2,liang2024osd,rajabi2025subtrackpp,rajabi2024grasstracklllm,refael2025sumo,apollo2024,grass2024,lotus2026}. They are motivated by empirical evidence that neural gradients have persistent low-dimensional structure \cite{thomas2026nosubspace,huang2026spectra,gradsubspace_survey2023}. However, a single subspace is a strong modeling choice: for example, a PINN layer may contain a slowly moving core induced by the PDE residual while also needing fast directions near transient boundary corrections; a graph model may begin with highly aligned neighborhood-aggregation gradients and lose that alignment as representations specialize; and a vision classifier may show low-rank structure that is stable enough for compression but not predictive of optimization gains.

This paper studies whether a \textbf{flag-manifold optimizer state} is a useful refinement of the single-subspace assumption. A flag is a nested chain $U_1\subset U_2\subset\cdots\subset U_K$, so it can represent a stable inner core and progressively larger outer directions. STRATA-AdamW applies this idea to optimizer state: it maintains multi-timescale covariance sketches, updates a flag by Riemannian steps, aligns blocks by Procrustes transport, and applies Adam-style moments in the incremental flag blocks. This construction is related to mathematical work on flag optimization \cite{ye2022flag,lai2022simpler,szwagier2025nested,jin2025online,funwithflags2024,xiao2026flagvarieties,ye2021flagtool,flagdecomp2025}, but the object being optimized is different: STRATA uses a flag inside a stochastic optimizer for changing neural-network gradients.

Our study is intentionally not written as a state-of-the-art claim. The original implementation exposed two numerical specification bugs that are easy to miss because they arise only after projecting gradients into low-dimensional flag blocks. First, transporting a mature first moment while resetting the second moment but sharing a single Adam bias-correction counter over-corrects the first moment after each flag update. Second, using the usual $10^{-8}$ Adam epsilon in near-zero projected blocks can divide by a vanishing second moment, producing huge block updates. Fixing these issues makes the optimizer much more stable, but it does not turn STRATA into a universal winner.

\begin{table*}[t]
\centering
\caption{\textbf{One STRATA-AdamW iteration.} The parameter branch uses the current flag; refreshed coordinates are used at the next iteration. Block alignment rotates first-moment coordinates but need not preserve ambient momentum exactly.}
\label{tab:method-update}
\begin{tabularx}{\textwidth}{@{}l>{\raggedright\arraybackslash}p{.31\textwidth}>{\raggedright\arraybackslash}X@{}}
\toprule
\textbf{Stage} & \textbf{Operation} & \textbf{State or output} \\
\midrule
Project & Decompose the gradient into orthogonal incremental flag blocks. & $Z_{k,t}=G_tU_t^{[k]}$, \quad $G_{\perp,t}=G_t(I-P_{K,t})$. \\
Adapt & Update block first and second moments with separate ages. & $\widehat M_{k,t}$ and $\widehat V_{k,t}$ use counters $c_k^M$ and $c_k^V$. \\
Lift & Lift the blockwise adaptive directions; add the normalized residual. & $D_t=\sum_k a_kH_{k,t}U_t^{[k]\top}+a_\perp G_{\perp,t}/(\|G_{\perp,t}\|_F/\sqrt{md}+\epsilon)$. \\
Update & Apply the direction and decoupled weight decay. & Unconstrained matrix parameter $W_{t+1}$. \\
\midrule
\rowcolor{oursbg}
Track & Update covariance sketches at every step. & $Y_{k,t}=\rho_kY_{k,t-1}+(1-\rho_k)G_t^\top(G_t\Omega_k)$. \\
\rowcolor{oursbg}
Refresh & Every $\tau$ steps, take a flag step, retract, and align blocks. & Procrustes rotations transport $M_k$; reset $V_k$ and $c_k^V$, retain mature $c_k^M$. \\
\bottomrule
\end{tabularx}
\end{table*}



The completed follow-ups identify targeted benefits in three of twelve citation-graph settings (Table~\ref{tab:positive-cases}), with every remaining setting retained in Appendix~\ref{app:full-empirical-record}. On PubMed/SAGE, independently validation-selected STRATA improves on both same-rank Single-Grassmann and AdamW, while its measured optimizer tensor state is smaller than AdamW's (Table~\ref{tab:pubmed-resources}). This does not establish a family-wide ranking: GaLore and Independent-Grassmann remain stronger in that cohort, several other graph and tabular settings favor the baselines, and the medium-scale image searches are still running. An approximately state-matched Single-Grassmann control wins at the default flag step; equal six-candidate flag-step/rate searches yield a smaller STRATA gain of 0.460 percentage points (Table~\ref{tab:memory-control}). We therefore retain budget and hyperparameter matching when testing the positive validation signal across a balanced architecture matrix and mechanism controls. Equalizing flag-layer coefficients barely changes the best archived Burgers or Poisson cells. In the new PubMed/SAGE grid, by contrast, all ten validation-selected interventions reduce test accuracy relative to full STRATA; these setting-dependent results support testing the combined design without identifying a universal causal mechanism. Our contributions are: (i) we identify and fix two flag-projection-specific numerical stability bugs; (ii) we introduce $\delta_{\mathrm{nest}}$ as a descriptive diagnostic of multi-timescale gradient nesting, while explicitly showing that it is not a performance predictor; (iii) we provide a fully proved conditional analysis connecting flag optimality, tracking, perturbation, nonconvex and PL convergence, and full-state generalization stability, with a transparent accounting of which quantities are measured in this project and which remain out of scope.
